\documentclass[a4paper,11pt]{article}
\usepackage[T2A]{fontenc}
\usepackage[utf8]{inputenc}
\usepackage[russian]{babel}
\usepackage{amsmath,amssymb,mathtools}
\usepackage{geometry}
\geometry{left=23mm,right=23mm,top=22mm,bottom=23mm}
\usepackage{microtype}
\usepackage{icomma}
\usepackage{tikz}
\usetikzlibrary{calc,arrows.meta,positioning,shapes.geometric}
\usepackage{pgfplots}
\pgfplotsset{compat=1.18}
\usepackage{booktabs,tabularx}
\usepackage{enumitem}
\usepackage{fancyhdr}
\usepackage{setspace}
\usepackage{needspace}
\usepackage{xcolor}
\usepackage{hyperref}
\definecolor{accent}{HTML}{008080}
\definecolor{darkaccent}{HTML}{004D4D}
\definecolor{textgray}{HTML}{333333}
\definecolor{lightline}{HTML}{B7CACA}
\hypersetup{colorlinks=true,linkcolor=darkaccent,pdfauthor={Лёвин Артём Александрович},pdftitle={Чек-лист: линейная функция и взаимное расположение прямых}}
\onehalfspacing
\setlength{\parindent}{0pt}
\setlength{\parskip}{5pt}
\setlist[itemize]{leftmargin=18pt,itemsep=2pt,topsep=3pt}
\setlist[enumerate]{leftmargin=23pt,itemsep=6pt,topsep=4pt}
\pagestyle{fancy}
\fancyhf{}
\fancyhead[L]{\small\color{textgray} Линейная функция и графики прямых}
\fancyhead[R]{\small\color{textgray} Кирилл Зиновьев}
\fancyfoot[C]{\small\thepage}
\renewcommand{\headrulewidth}{0.4pt}
\renewcommand{\headrule}{\hbox to\headwidth{\color{accent}\leaders\hrule height \headrulewidth\hfill}}
\setlength{\headheight}{23pt}
\newcommand{\topic}[1]{\vspace{7pt}{\color{darkaccent}\large\bfseries #1}\par\vspace{-4pt}{\color{accent}\rule{\linewidth}{0.45pt}}\par\vspace{3pt}}
\newcommand{\minorheading}[1]{\needspace{4\baselineskip}\vspace{2pt}{\color{darkaccent}\bfseries #1}\par}
\newcommand{\keyidea}[1]{{\color{darkaccent}\bfseries Ключевая идея.} #1}
\newcommand{\citem}[1]{\(\square\)\enspace #1\par}
\begin{document}
\thispagestyle{empty}
\begin{center}
\vspace*{24mm}
{\color{darkaccent}\large\MakeUppercase{Математика}}\par
\vspace{25mm}
{\fontsize{27}{32}\selectfont Чек-лист}\par
\vspace{17mm}
{\color{accent}\fontsize{19}{25}\selectfont Линейная функция:\par
угловой коэффициент и\par
взаимное расположение прямых}\par
\vspace{19mm}
{\large Координатная плоскость и графики функций}\par
\vspace{33mm}
{\large Кирилл Зиновьев}\par
\vspace{12mm}
{\large 15 ноября 2026 г.}\par
\vfill
{\small Лёвин Артём Александрович\par эксклюзивно для Кирилла Зиновьева}\par
\vspace{14mm}
\begin{tikzpicture}[x=1cm,y=1cm]
\draw[accent,line width=.9pt] (0,0) .. controls (1.7,.65) and (3.1,-.52) .. (4.7,0) .. controls (6.25,.55) and (7.9,-.55) .. (9.5,0);
\end{tikzpicture}
\end{center}
\clearpage

\topic{1. Что означают коэффициенты в формуле прямой}
\keyidea{График линейной функции имеет вид прямой.}
\[
\boxed{y=kx+b}
\]
Здесь \(x\) --- независимая переменная, \(y\) --- значение функции, \(k\) --- \textbf{угловой коэффициент}, \(b\) --- \textbf{свободный член}.
\begin{itemize}
\item \(k\) отвечает за направление и наклон прямой;
\item \(b\) показывает, где прямая пересекает ось \(Oy\): в точке \((0;b)\);
\item коэффициент может быть отрицательным, положительным или равным нулю.
\end{itemize}
\minorheading{Как распознать \(k\)}
\begin{center}
\begin{tabular}{@{}lcc@{}}
\toprule
\textbf{Функция} & \(k\) & \(b\)\\
\midrule
\(y=3x+4\) & \(3\) & \(4\)\\
\(y=-x+4\) & \(-1\) & \(4\)\\
\(y=x-5\) & \(1\) & \(-5\)\\
\(y=7\) & \(0\) & \(7\)\\
\bottomrule
\end{tabular}
\end{center}
Обязательно помните: \(-x=(-1)x\), а \(x=1\cdot x\). Число перед \(x\) может быть не записано явно.

\topic{2. Как угловой коэффициент влияет на график}
\begin{itemize}
\item \(k>0\): прямая поднимается слева направо;
\item \(k<0\): прямая опускается слева направо;
\item \(k=0\): прямая горизонтальна, \(y=b\).
\end{itemize}
\[
\underbrace{y=0\cdot x+4}_{k=0}\quad\Longrightarrow\quad \boxed{y=4}.
\]
\keyidea{Горизонтальная прямая параллельна оси абсцисс \(Ox\), если \(b\ne0\).}
При \(b=0\) прямая \(y=0\) совпадает с осью \(Ox\).

\topic{3. Параллельные прямые}
Для \(y=k_1x+b_1\) и \(y=k_2x+b_2\) угловые коэффициенты равны при одинаковом направлении прямых:
\[
\boxed{k_1=k_2}.
\]
Если при этом \(b_1\ne b_2\), прямые \textbf{различны и параллельны}. Если и свободные члены равны, прямые совпадают.

\minorheading{Пример 1. Найти коэффициент по условию параллельности}
При каком \(k\) график \(y=kx+4\) параллелен прямой \(y=-x\)?
\[
 y=-x=(-1)x+0\quad\Longrightarrow\quad \boxed{k=-1}.
\]
Получается \(y=-x+4\). Одинаковые угловые коэффициенты равны \(-1\), а свободные члены различны: \(4\) и \(0\).

\begin{center}
\begin{tikzpicture}
\begin{axis}[
 width=.73\linewidth,height=6.1cm,axis lines=middle,
 xmin=-2.5,xmax=2.6,ymin=-5,ymax=8,
 xtick={-2,-1,0,1,2},ytick={-4,-2,0,2,4,6},
 xlabel={$x$},ylabel={$y$},
 grid=major,grid style={lightline!45},
 axis line style={->},tick label style={font=\scriptsize},
 label style={font=\small},clip=false]
\addplot[accent,very thick,domain=-2.5:2.6,samples=2]{2*x+2};
\addplot[darkaccent,very thick,densely dashed,domain=-2.5:2.6,samples=2]{2*x-2};
\node[anchor=west,font=\small,fill=white,inner sep=1pt,text=accent] at (axis cs:.45,5.4) {$y=2x+2$};
\node[anchor=west,font=\small,fill=white,inner sep=1pt,text=darkaccent] at (axis cs:.5,-2.3) {$y=2x-2$};
\end{axis}
\end{tikzpicture}
\par\small\color{textgray} Прямые имеют одинаковый наклон \(k=2\), но разные свободные члены.
\end{center}

\topic{4. Оси координат и точки пересечения}
\textbf{Ось абсцисс} \(Ox\) --- горизонтальная ось; \textbf{ось ординат} \(Oy\) --- вертикальная ось.
\begin{center}
\begin{tabular}{@{}lcc@{}}
\toprule
\textbf{Ось} & \textbf{Координата} & \textbf{Вид точки}\\
\midrule
\(Ox\) & \(y=0\) & \((x;0)\)\\
\(Oy\) & \(x=0\) & \((0;y)\)\\
\bottomrule
\end{tabular}
\end{center}
\minorheading{Пример 2. Пересечение с осью \(Ox\)}
Найдите точку пересечения прямой \(y=0{,}4x-12\) с осью абсцисс.
\[
0=0{,}4x-12\quad\Longrightarrow\quad
0{,}4x=12\quad\Longrightarrow\quad x=30.
\]
\textbf{Ответ:} \(A(30;0)\). Именно пара координат задаёт искомую точку.

\minorheading{Когда прямая \(y=kx+b\) не пересекает \(Ox\)?}
\[
\boxed{k=0,\quad b\ne0}.
\]
Например, \(y=4\) целиком лежит выше оси \(Ox\), поэтому общих точек с нею нет.

\topic{5. Как найти \(k\), если известна точка прямой}
Если прямая \(y=kx+b\) проходит через точку \(T(x_0;y_0)\), её координаты \textbf{обязательно удовлетворяют} уравнению:
\[
\boxed{y_0=kx_0+b}.
\]
\minorheading{Пример 3. Точка на оси абсцисс}
При каком \(k\) прямая \(y=kx+4\) проходит через \(T(3;0)\)?
\[
0=3k+4\quad\Longrightarrow\quad 3k=-4
\quad\Longrightarrow\quad\boxed{k=-\frac43}.
\]
\minorheading{Пример 4. Произвольная точка}
Прямая \(y=kx-1\) проходит через точку \(T(-1;-4)\). Тогда
\[
-4=k\cdot(-1)-1=-k-1
\quad\Longrightarrow\quad -3=-k
\quad\Longrightarrow\quad \boxed{k=3}.
\]
\textbf{Проверка:} \(3\cdot(-1)-1=-4\), точка действительно принадлежит прямой.

Если \(x_0=0\), значение \(k\) таким способом обычно определить нельзя: необходимо сначала проверить равенство \(y_0=b\).

\topic{6. Как найти точку пересечения двух прямых}
В общей точке пересечения обе функции имеют \textbf{одинаковые \(x\) и \(y\)}. Поэтому их правые части нужно приравнять.

\minorheading{Пример 5. Задача из двух этапов}
График \(y=kx+4\) проходит через точку пересечения прямых
\[
y=12-x,\qquad y=x+4.
\]
\textbf{Шаг 1. Найдём пересечение заданных прямых:}
\[
12-x=x+4\quad\Longrightarrow\quad 8=2x\quad\Longrightarrow\quad x=4.
\]
\[
y=4+4=8,\qquad T(4;8).
\]
\textbf{Шаг 2. Подставим точку в искомую функцию:}
\[
8=4k+4\quad\Longrightarrow\quad 4k=4
\quad\Longrightarrow\quad\boxed{k=1}.
\]
\textbf{Проверка:} при \(k=1\) получаем \(y=x+4\); точка \((4;8)\) лежит на этой прямой. В данном примере искомый график совпадает с одной из двух исходных прямых, что допустимо по условию.

\clearpage
\thispagestyle{fancy}
\begin{center}
{\color{darkaccent}\Large\bfseries Алгоритм: прямая через точку пересечения}\par
\vspace{10mm}
\end{center}
\tikzset{
 flowstart/.style={ellipse,draw=accent,very thick,minimum width=8cm,minimum height=12mm,align=center,font=\normalsize},
 flowstep/.style={rectangle,rounded corners=2mm,draw=accent,line width=.8pt,minimum width=11.7cm,text width=10.8cm,minimum height=17mm,align=center,font=\normalsize,inner sep=4pt},
 flowarrow/.style={-{Latex[length=2.9mm,width=1.8mm]},line width=.85pt,draw=darkaccent}
}
\begin{center}
\begin{tikzpicture}[node distance=10mm]
\node[flowstart] (s) {Даны две прямые и функция \(y=kx+b\)};
\node[flowstep,below=of s] (a) {Приравняйте правые части уравнений двух прямых};
\node[flowstep,below=of a] (b) {Решите полученное уравнение и найдите \(x_0\)};
\node[flowstep,below=of b] (c) {Подставьте \(x_0\) в любое исходное уравнение; найдите \(y_0\)};
\node[flowstep,below=of c] (d) {Подставьте \((x_0;y_0)\) в \(y=kx+b\)};
\node[flowstep,below=of d] (e) {Решите уравнение относительно \(k\)};
\node[flowstep,below=of e] (f) {Проверьте, что найденная прямая проходит через точку};
\node[flowstart,below=of f] (g) {Ответ: найденный коэффициент \(k\)};
\foreach \u/\v in {s/a,a/b,b/c,c/d,d/e,e/f,f/g}
  \draw[flowarrow] (\u.south) -- (\v.north);
\end{tikzpicture}
\end{center}
\clearpage

\topic{7. Ошибки, которые важно предупреждать}
\begin{enumerate}[label=\arabic*.]
\item \textbf{Путают коэффициенты.} В \(y=kx+b\) наклон задаёт \(k\), а пересечение с \(Oy\) --- \(b\).
\item \textbf{Забывают знак перед \(x\).} У функции \(y=-x\) коэффициент \(k=-1\).
\item \textbf{Приравнивают свободные члены при проверке параллельности.} Для одинакового направления сравнивайте \(k\); если \(b\) также равны, прямые совпадают.
\item \textbf{Путают оси.} На \(Ox\) всегда \(y=0\); на \(Oy\) всегда \(x=0\).
\item \textbf{Считывают точку в обратном порядке.} \(T(3;0)\) означает \(x=3\), \(y=0\).
\item \textbf{Меняют знак при перестановке всего равенства.} Из \(0=3k+4\) можно написать \(3k+4=0\) без изменения знаков. При переносе \emph{отдельного слагаемого} знак меняется: \(3k=-4\).
\item \textbf{Не заканчивают задачу.} Если требуется точка, запишите обе координаты; если требуется \(k\), дайте значение коэффициента.
\end{enumerate}

\topic{8. Мини-чек-лист самостоятельной проверки}
\citem{Я различаю угловой коэффициент \(k\) и свободный член \(b\).}
\citem{Я помню условие параллельности и случай совпадения прямых.}
\citem{Я могу объяснить, почему при \(k=0\) прямая горизонтальна.}
\citem{Я знаю, какая координата равна нулю на каждой оси.}
\citem{Я умею подставлять \((x_0;y_0)\) в уравнение функции.}
\citem{Я нахожу пересечение двух прямых, приравнивая их формулы.}
\citem{Я проверяю найденный коэффициент подстановкой.}

\minorheading{Краткая памятка}
\renewcommand{\arraystretch}{1.35}
\begin{tabularx}{\linewidth}{@{}p{.47\linewidth}X@{}}
\toprule
\textbf{Условие задачи} & \textbf{Первое действие} \\
\midrule
Прямая параллельна другой прямой & Приравнять угловые коэффициенты \\
Прямая горизонтальна & Положить \(k=0\) \\
Прямая проходит через \((x_0;y_0)\) & Подставить обе координаты \\
Нужно пересечение с \(Ox\) & Положить \(y=0\) \\
Нужно пересечение с \(Oy\) & Положить \(x=0\) \\
Нужна общая точка двух прямых & Приравнять правые части \\
\bottomrule
\end{tabularx}

\clearpage
\topic{Тренировочный блок --- путь А}
\textbf{10 заданий.} Выполняйте их по порядку: от распознавания коэффициентов до комбинированных задач. В заданиях, где требуется найти \(k\), обязательно проверьте ответ подстановкой.
\vspace{3mm}
{\setstretch{1.18}
\begin{enumerate}[label=\textbf{\arabic*.},itemsep=8pt]
\item Укажите угловой коэффициент и свободный член функции \(y=-3x+7\).
\item При каком \(k\) прямая \(y=kx-5\) параллельна \(y=4x+1\)?
\item При каком \(k\) прямая \(y=kx+6\) параллельна \(y=-x\)?
\item При каком \(k\) график \(y=kx-9\) параллелен оси \(Ox\)?
\item При каком \(k\) прямая \(y=kx+2\) проходит через точку \(A(2;8)\)?
\item При каком \(k\) прямая \(y=kx-7\) проходит через точку \(B(-2;-1)\)?
\item Найдите координаты точки пересечения графика \(y=0{,}6x-9\) с осью \(Ox\).
\item При каком \(k\) график \(y=kx-10\) пересекает ось \(Ox\) в точке с абсциссой \(5\)?
\item Найдите \(k\), если график \(y=kx+1\) проходит через точку пересечения прямых \(y=9-x\) и \(y=2x+3\).
\item Найдите \(k\), если график \(y=kx-3\) проходит через точку пересечения прямых \(y=10-2x\) и \(y=x+1\).
\end{enumerate}
}
\vspace{4mm}
\textbf{Проверка готовности:} после решения попробуйте словами объяснить, какая идея использована в каждом задании: равенство наклонов, подстановка координат или поиск точки пересечения.

\clearpage
\topic{Ответы к тренировочному блоку}
\vspace{6mm}
\begin{center}
\renewcommand{\arraystretch}{1.35}
\begin{tabular}{@{}c p{11.5cm}@{}}
\toprule
\textbf{№} & \textbf{Ответ}\\
\midrule
1 & \(k=-3,\quad b=7\)\\
2 & \(k=4\)\\
3 & \(k=-1\)\\
4 & \(k=0\)\\
5 & \(k=3\)\\
6 & \(k=-3\)\\
7 & \((15;0)\)\\
8 & \(k=2\)\\
9 & \(k=3\)\\
10 & \(k=\dfrac{7}{3}\)\\
\bottomrule
\end{tabular}
\end{center}
\vspace{4mm}
\topic{Итоговая памятка}
\begin{center}
\renewcommand{\arraystretch}{1.25}
\begin{tabularx}{\linewidth}{@{}X X@{}}
\toprule
\textbf{Вопрос} & \textbf{Опорная формула или условие}\\
\midrule
Когда две различные прямые параллельны? & \(k_1=k_2\), \(b_1\ne b_2\)\\
Когда прямая горизонтальна? & \(k=0\)\\
Что выполняется в точке на \(Ox\)? & \(y=0\)\\
Что выполняется в точке на \(Oy\)? & \(x=0\)\\
Как найти \(k\) по точке? & \(y_0=kx_0+b\)\\
Как искать общую точку двух прямых? & Приравнять значения \(y\)\\
\bottomrule
\end{tabularx}
\end{center}
\end{document}
